% !TEX program = pdflatex
\documentclass[11pt]{article}

\usepackage[T1]{fontenc}
\usepackage{lmodern}
\usepackage{microtype}
\usepackage[a4paper,margin=28mm,headheight=14pt]{geometry}
\usepackage{amsmath,amssymb,amsthm,mathtools}
\usepackage{enumitem}
\usepackage{array,booktabs,longtable}
\usepackage{xcolor}
\usepackage{xurl}
\usepackage{hyperref}
\usepackage{fancyhdr}

\definecolor{linkblue}{RGB}{24,78,119}
\definecolor{citegreen}{RGB}{36,104,71}
\hypersetup{
  colorlinks=true,
  linkcolor=linkblue,
  citecolor=citegreen,
  urlcolor=linkblue,
  pdftitle={A Dirichlet Polya Inequality for Three-Dimensional Spherical Shells},
  pdfsubject={Purely analytic main proof manuscript},
  pdfkeywords={Polya conjecture, Dirichlet Laplacian, spherical shell, eigenvalue counting}
}

\pagestyle{fancy}
\fancyhf{}
\fancyhead[L]{\small Dirichlet P\'olya inequality for spherical shells}
\fancyhead[R]{\small Purely analytic proof}
\fancyfoot[C]{\thepage}
\fancypagestyle{plain}{%
  \fancyhf{}%
  \fancyhead[L]{\small Dirichlet P\'olya inequality for spherical shells}%
  \fancyhead[R]{\small Purely analytic proof}%
  \fancyfoot[C]{\thepage}%
}
\fancypagestyle{titlepage}{%
  \fancyhf{}%
  \fancyfoot[C]{\thepage}%
  \renewcommand{\headrulewidth}{0pt}%
}

\numberwithin{equation}{section}
\allowdisplaybreaks[2]
\setlist{itemsep=0.25em,topsep=0.45em}
\setlength{\emergencystretch}{3em}
\newtheorem{theorem}{Theorem}[section]
\newtheorem{proposition}[theorem]{Proposition}
\newtheorem{lemma}[theorem]{Lemma}
\newtheorem{corollary}[theorem]{Corollary}
\theoremstyle{definition}
\newtheorem{definition}[theorem]{Definition}
\theoremstyle{remark}
\newtheorem{remark}[theorem]{Remark}

\newcommand{\R}{\mathbb R}
\newcommand{\N}{\mathbb N}
\newcommand{\Nzero}{\mathbb N_0}
\DeclareMathOperator{\arsinh}{arsinh}

\title{A Dirichlet P\'olya Inequality for\\Three-Dimensional Spherical Shells\\[0.4em]
\large Purely analytic main proof manuscript}
\author{Anonymous for peer review}
\date{17 July 2026}

\begin{document}

\maketitle
\thispagestyle{titlepage}

\begin{quote}
\small
\textbf{Scope and status.}
This manuscript proves a theorem for the class of genuine three-dimensional
spherical shells.  It does not claim the general P\'olya conjecture,
literature novelty, or priority.  All shell-specific analytic reductions,
regional arguments, boundary assignments, exact constants, and all finite
exact-rational decisions are stated in this paper and its accompanying
purely analytic supplement.  Executable replay sources and interval
certificates are retained only as optional regression checks and are not
premises of the proof.  The imported background and analytic results
comprise explicitly identified standard Sturm--Liouville and Bessel facts and
published phase, Bessel-zero, and one-dimensional floor-sum inequalities,
each quoted with its precise scope.
\end{quote}

\begin{abstract}
For \(0<r<R\), let
\[
  A_{r,R}=\{x\in\R^3:r<|x|<R\}.
\]
We prove, with the strict Dirichlet counting convention, that
\[
  N_D(A_{r,R},\Lambda)
  \le L_3|A_{r,R}|\Lambda^{3/2},
  \qquad
  L_3=\frac1{6\pi^2},
  \qquad \Lambda\ge0.
\]
The proof is purely analytic.  It separates variables, converts the radial
spectrum into a strict multiplicity-weighted phase count, and combines an
exact retained-remainder identity with a ratio-sharp angular block estimate,
a tangent-envelope radial minorant, a universal ball quarter-layer, and an
elementary scalar convexity argument.  The low- and middle-ratio proof starts
at the first possible radial frequency; no finite Bessel-zero staircase,
event-state theorem, floating-point conclusion, or interval-arithmetic
conclusion is a proof premise.  We also give an elementary proof that no genuine spherical
shell tiles \(\R^3\) by congruent rigid-motion copies with disjoint interiors,
under either exact or almost-everywhere coverage.
\end{abstract}

\tableofcontents
\clearpage

\section{Main results and conventions}

The eigenvalues of the Dirichlet Laplacian on a bounded domain \(\Omega\) are
listed with multiplicity as
\[
  0<\lambda_1^D(\Omega)\le\lambda_2^D(\Omega)\le\cdots.
\]
Throughout we use the strict counting function
\begin{equation}
  N_D(\Omega,\Lambda)
  :=\#\{j:\lambda_j^D(\Omega)<\Lambda\}.
  \label{eq:strict-count}
\end{equation}
This convention matters at every spectral wall.  We write
\(\N=\{1,2,3,\ldots\}\) and \(\Nzero=\{0,1,2,\ldots\}\).

\begin{theorem}[Spectral inequality]
\label{thm:spectral}
For every \(0<r<R\) and every \(\Lambda\ge0\),
\begin{equation}
  \boxed{
  N_D(A_{r,R},\Lambda)
  \le \frac{2}{9\pi}(R^3-r^3)\Lambda^{3/2}
  =L_3|A_{r,R}|\Lambda^{3/2},
  \qquad L_3=\frac1{6\pi^2}.}
  \label{eq:main-theorem}
\end{equation}
\end{theorem}

The standard non-strict formulation follows without changing the constant.
Indeed, if
\[
 N_D^{\leq}(\Omega,\Lambda)
 :=\#\{j:\lambda_j^D(\Omega)\leq\Lambda\},
\]
then discreteness of the Dirichlet spectrum gives
\begin{equation}
 N_D^{\leq}(\Omega,\Lambda)
 =\lim_{\varepsilon\downarrow0}
 N_D(\Omega,\Lambda+\varepsilon).
 \label{eq:strict-to-nonstrict}
\end{equation}
Applying Theorem~\ref{thm:spectral} at
\(\Lambda+\varepsilon\) and passing to the limit proves
\eqref{eq:main-theorem} with \(N_D^{\leq}\) as well.  Thus the strict
convention used for the wall bookkeeping is equivalent here to the usual
non-strict P\'olya inequality.

\begin{theorem}[Non-tiling]
\label{thm:nontiling}
For every \(0<r<R\), no family of congruent rigid-motion copies of
\(A_{r,R}\) has pairwise-disjoint interiors and covers \(\R^3\) exactly or
up to a three-dimensional Lebesgue-null set.  The same conclusion holds when
coverage is formulated using closed shells.
\end{theorem}

The two theorems concern literally the same complete shell family.
Theorem~\ref{thm:nontiling} is logically independent of the spectral proof.

\subsection*{Proof map}

For \(K>0\), the spectral argument has one global comparison, a no-mode
region, and two high-frequency ratio owners.  After scaling to
\(A_{\rho,1}\), separation of variables and the uniform Bessel-phase
enclosure give a strict multiplicity-weighted lattice proxy \(P\) with
\[
 N_D(A_{\rho,1},K^2)\leq P.
\]
In the original action variables of Part~VII of the supplement, let \(R\) be
the inverse action, \(F=R^2\), and let \(M(r)\) count the half-integers
strictly below \(r\).  The inverse action gives the exact layer-cake picture
\[
 P=\sum_{\substack{n\geq1\\n-1/4<T}}M\bigl(R(n-1/4)\bigr)^2,
 \qquad
 W=\int_0^T R(t)^2\,dt,
\]
whose scaled version is Lemma~\ref{thin:lem-layer-cake}.  Thus the master
comparison is
\begin{equation}
 W-P=\mathcal D_{\rm rad}-\mathcal E_{\rm ang}:
 \quad\hbox{radial quadrature reserve pays angular rounding.}
 \label{eq:proof-map-master-defect}
\end{equation}
The exact definitions and the retained positive remainders in
\eqref{eq:proof-map-master-defect} are given in the accompanying
\emph{Purely Analytic Supplement}.

\begin{samepage}
The parameter space is closed by the following disjoint owners.
\begin{center}
\begin{tabular}{@{}lll@{}}
\toprule
frequency range&ratio range&owner\\
\midrule
\(0\leq K\leq\pi/(1-\rho)\)&\(0<\rho<1\)&strict no-mode bound\\
\(K>\pi/(1-\rho)\)&\(0<\rho<39/50\)&global defect theorem\\
\(K>\pi/(1-\rho)\)&\(39/50\leq\rho<1\)&optical theorem\\
\bottomrule
\end{tabular}
\end{center}
\end{samepage}
Their union proves the unit-shell inequality; dilation proves the result for
every \(A_{r,R}\).  The former finite staircase, 38-state theorem, and
611-row allocation are retained in the repository only as historical audit
material; they have no implication into the proof below.

% Self-contained manuscript fragment: separated spectrum, phase reduction,
% weighted lattice reduction, and the small-hole endpoint.
% Assumes amsmath, amssymb, amsthm, mathtools and theorem environments
% theorem, proposition, lemma.

\section{Separated spectrum and phase reduction}
\label{sec:sc-spectrum-small-hole}

Throughout this section
\[
  A_{\rho,1}=\{x\in\mathbb R^3:\rho<|x|<1\},
  \qquad 0<\rho<1,
\]
and the Dirichlet counting function is strict:
\[
  N_D(A_{\rho,1},K^2)
  :=\#\{j:\lambda_j^D(A_{\rho,1})<K^2\}.
\]
For \(x>0\) we use the strict positive-integer count
\begin{equation}
  [x]_<:=\#\{n\in\mathbb N:n<x\}
  =\max\{0,\lceil x\rceil-1\}.
  \label{eq:sc-strict-bracket}
\end{equation}
Thus \([m]_< =m-1\) for \(m\in\mathbb N\).  This convention is retained
at every eigenfrequency, phase level, and half-integer angular wall.

\subsection{The exact separated spectrum}

Let \(H_\rho=-\Delta^D_{A_{\rho,1}}\) be the Friedrichs operator associated
with
\[
  \mathfrak q[f]=\int_{A_{\rho,1}}|\nabla f|^2\,dx,
  \qquad D(\mathfrak q)=H_0^1(A_{\rho,1}).
\]
Choose an orthonormal spherical-harmonic basis satisfying
\[
  -\Delta_{\mathbb S^2}Y_{\ell m}
  =\ell(\ell+1)Y_{\ell m},
  \qquad \ell\in\mathbb N_0,
  \quad -\ell\le m\le\ell.
\]

\begin{proposition}[Unitary decomposition and completeness]
\label{prop:sc-separated-spectrum}
For
\[
  L_\ell=-\frac{d^2}{dr^2}+\frac{\ell(\ell+1)}{r^2},
  \qquad
  D(L_\ell)=H^2(\rho,1)\cap H_0^1(\rho,1),
\]
one has the unitary equivalence
\begin{equation}
  H_\rho\simeq
  \bigoplus_{\ell=0}^{\infty}
  \bigoplus_{m=-\ell}^{\ell}L_\ell.
  \label{eq:sc-operator-sum}
\end{equation}
The exact form domain on the right is
\begin{equation}
  \left\{(u_{\ell m}):
  \begin{array}{l}
  u_{\ell m}\in H_0^1(\rho,1),\\
  \displaystyle
  \sum_{\ell,m}\left(\|u_{\ell m}\|_2^2
  +\mathfrak a_\ell[u_{\ell m}]\right)<\infty
  \end{array}\right\},
  \label{eq:sc-global-form-domain}
\end{equation}
where
\[
  \mathfrak a_\ell[u]
  =\int_\rho^1\left(|u'|^2
  +\frac{\ell(\ell+1)}{r^2}|u|^2\right)dr.
\]
The exact operator domain is
\begin{equation}
  \left\{(u_{\ell m}):
  \begin{array}{l}
  u_{\ell m}\in D(L_\ell),\\
  \displaystyle
  \sum_{\ell,m}\left(\|u_{\ell m}\|_2^2
  +\|L_\ell u_{\ell m}\|_2^2\right)<\infty
  \end{array}\right\}.
  \label{eq:sc-global-operator-domain}
\end{equation}
Each \(L_\ell\) has positive, simple eigenvalues tending to infinity and a
complete orthonormal eigenbasis.  The full operator has compact resolvent.
\end{proposition}

\begin{proof}
Write \(x=r\omega\), so \(dx=r^2\,dr\,d\omega\), and set
\[
  R_{\ell m}(r)=\int_{\mathbb S^2}
  f(r,\omega)\overline{Y_{\ell m}(\omega)}\,d\omega,
  \qquad u_{\ell m}(r)=rR_{\ell m}(r).
\]
Fubini, Tonelli, and Parseval on \(\mathbb S^2\) give
\[
  \|f\|_{L^2(A_{\rho,1})}^2
  =\sum_{\ell,m}\int_\rho^1|R_{\ell m}(r)|^2r^2\,dr
  =\sum_{\ell,m}\int_\rho^1|u_{\ell m}(r)|^2\,dr.
\]
Conversely, if \((u_{\ell m})\) is square summable, the finite sums
\[
  f_L(r,\omega)=\frac1r
  \sum_{\ell\le L}\sum_{m=-\ell}^{\ell}
  u_{\ell m}(r)Y_{\ell m}(\omega)
\]
are Cauchy in \(L^2(A_{\rho,1})\), and their limit maps back to the given
family.  Hence the coefficient map is unitary onto the Hilbert direct sum.
Because \(r,1/r\in W^{1,\infty}(\rho,1)\), multiplication by either function
is bounded on \(H^1(\rho,1)\) and preserves the two zero endpoint traces.

For a finite smooth harmonic expansion, polar coordinates and angular
orthogonality yield
\[
  \mathfrak q[f]=\sum_{\ell,m}\int_\rho^1
  \left(r^2|R_{\ell m}'|^2
  +\ell(\ell+1)|R_{\ell m}|^2\right)dr.
\]
If \(u=rR\), then
\[
  r^2|R'|^2
  =\left|u'-\frac ur\right|^2
  =|u'|^2-\frac{(|u|^2)'}r+\frac{|u|^2}{r^2}.
\]
Since
\[
  \frac{(|u|^2)'}r
  =\left(\frac{|u|^2}{r}\right)'
  +\frac{|u|^2}{r^2}
\]
and \(u(\rho)=u(1)=0\), integration gives
\[
  \int_\rho^1r^2|R'|^2\,dr
  =\int_\rho^1|u'|^2\,dr.
\]
The angular term becomes
\(\ell(\ell+1)|u|^2/r^2\).  Closing finite harmonic expansions in the
form norm proves \eqref{eq:sc-global-form-domain}: one inclusion follows
from componentwise form convergence and Fatou's lemma, while the reverse
inclusion follows by first truncating the square-summable family in
\((\ell,m)\) and then approximating its finitely many components by
\(C_c^\infty(\rho,1)\).  The uniqueness theorem for operators associated
with closed semibounded forms gives \eqref{eq:sc-operator-sum}.

For fixed \(\ell\), the potential \(\ell(\ell+1)/r^2\) is continuous and
bounded on \([\rho,1]\).  The regular one-dimensional Dirichlet form
therefore has associated operator \(L_\ell\) with the displayed
\(H^2\cap H_0^1\) domain.  The standard domain formula for an orthogonal
operator sum gives \eqref{eq:sc-global-operator-domain}.

The embedding \(H_0^1(\rho,1)\hookrightarrow L^2(\rho,1)\) is compact, so
every block has compact resolvent.  To pass to the infinite direct sum, note
that \(r\le1\) implies
\[
  L_\ell\ge\ell(\ell+1),
  \qquad
  \|(L_\ell+1)^{-1}\|
  \le\frac1{1+\ell(\ell+1)}.
\]
The resolvent of the full operator is consequently the operator-norm limit
of its finite-\(\ell\) truncations, each of which is compact.  Hence the full
resolvent is compact.

Regular Dirichlet Sturm--Liouville theory gives a complete sequence of
simple eigenvalues in every block.  They are strictly positive because
\[
  \mathfrak a_\ell[u]\ge\int_\rho^1|u'|^2\,dr
  \ge\frac{\pi^2}{(1-\rho)^2}\|u\|_2^2.
\]
Simplicity also follows directly: two solutions of the same second-order
equation that vanish at \(r=\rho\) have proportional initial derivatives
and hence are proportional by uniqueness.  If \(u_{\ell,n}\) is a complete
radial eigenbasis, then
\[
  \frac{u_{\ell,n}(r)}rY_{\ell m}(\omega)
\]
over all \((\ell,m,n)\) is exactly the inverse image of a complete
orthonormal coordinate basis of the Hilbert direct sum.  This proves global
completeness.
\end{proof}

\begin{proposition}[Bessel determinant and strict phase enumeration]
\label{prop:sc-bessel-phase-count}
Put \(\nu=\ell+\tfrac12\).  The positive eigenfrequencies in the
\(\ell\)-th radial block are precisely the positive zeros of
\begin{equation}
  D_{\nu,\rho}(k)
  :=J_\nu(\rho k)Y_\nu(k)-Y_\nu(\rho k)J_\nu(k).
  \label{eq:sc-bessel-determinant}
\end{equation}
Every such zero is simple.  With the continuous phase branch
\begin{equation}
  J_\nu(x)+iY_\nu(x)=M_\nu(x)e^{i\theta_\nu(x)},
  \qquad M_\nu(x)>0,
  \qquad \theta_\nu(0+)=-\frac\pi2,
  \label{eq:sc-phase-branch}
\end{equation}
and
\[
  \Psi_{\nu,\rho}(K)=\theta_\nu(K)-\theta_\nu(\rho K),
  \qquad
  \gamma_{\rho,K}(\nu)=\frac{\Psi_{\nu,\rho}(K)}\pi,
\]
the exact strict count is
\begin{equation}
  N_D(A_{\rho,1},K^2)
  =\sum_{\ell=0}^{\infty}(2\ell+1)
  [\gamma_{\rho,K}(\ell+\tfrac12)]_<.
  \label{eq:sc-exact-phase-count}
\end{equation}
The sum is finite.  If \(K\) is a radial eigenfrequency in one or several
angular channels, the root at \(K\) is excluded in each such channel and the
excluded angular multiplicities add.
\end{proposition}

\begin{proof}
For \(k>0\), the physical radial equation has the general solution
\[
  R(r)=r^{-1/2}
  \bigl(c_1J_\nu(kr)+c_2Y_\nu(kr)\bigr).
\]
The two Dirichlet conditions admit a nonzero coefficient vector exactly when
\[
  \det\begin{pmatrix}
  J_\nu(\rho k)&Y_\nu(\rho k)\\
  J_\nu(k)&Y_\nu(k)
  \end{pmatrix}=D_{\nu,\rho}(k)=0.
\]
The Wronskian
\begin{equation}
  J_\nu(x)Y_\nu'(x)-J_\nu'(x)Y_\nu(x)=\frac2{\pi x}
  \label{eq:sc-bessel-wronskian}
\end{equation}
shows that the two entries in either row cannot vanish simultaneously.
Thus the determinant criterion remains exact even if the \(J_\nu\) values,
or the \(Y_\nu\) values, happen to vanish at both endpoints.

For root simplicity define the left-normalized solution
\begin{equation}
  s_\ell(r,k)=\frac{\pi\sqrt\rho}{2}\sqrt r
  \left[J_\nu(\rho k)Y_\nu(kr)
  -Y_\nu(\rho k)J_\nu(kr)\right].
  \label{eq:sc-left-normalized}
\end{equation}
Equation \eqref{eq:sc-bessel-wronskian} gives
\[
  s_\ell(\rho,k)=0,
  \qquad \partial_rs_\ell(\rho,k)=1,
  \qquad
  s_\ell(1,k)=\frac{\pi\sqrt\rho}{2}D_{\nu,\rho}(k).
\]
Let \(v=\partial_ks_\ell\).  Differentiating
\((L_\ell-k^2)s_\ell=0\) gives
\[
  (L_\ell-k^2)v=2ks_\ell,
  \qquad v(\rho,k)=\partial_rv(\rho,k)=0.
\]
At a positive determinant root, Lagrange's identity yields
\begin{equation}
  \partial_ks_\ell(1,k)\,\partial_rs_\ell(1,k)
  =2k\int_\rho^1s_\ell(r,k)^2\,dr>0.
  \label{eq:sc-root-simplicity}
\end{equation}
Thus \(\partial_ks_\ell(1,k)\ne0\), and the zero of the determinant is
simple.

At zero energy the left-normalized solution is
\[
  s_\ell(r,0)=
  \frac{\rho^{-\ell}r^{\ell+1}
  -\rho^{\ell+1}r^{-\ell}}{2\ell+1},
\]
so
\[
  s_\ell(1,0)=
  \frac{\rho^{-\ell}-\rho^{\ell+1}}{2\ell+1}>0.
\]
Equivalently, the standard small-argument expansions give
\begin{equation}
  \lim_{k\downarrow0}D_{\nu,\rho}(k)
  =\frac{\rho^{-\nu}-\rho^\nu}{\pi\nu}>0.
  \label{eq:sc-zero-determinant-limit}
\end{equation}
Hence the limiting phase level zero is not an eigenfrequency.

From \eqref{eq:sc-phase-branch},
\[
  J_\nu=M_\nu\cos\theta_\nu,
  \qquad Y_\nu=M_\nu\sin\theta_\nu,
\]
and direct substitution gives the sign-sensitive factorization
\begin{equation}
  D_{\nu,\rho}(k)
  =M_\nu(\rho k)M_\nu(k)
  \sin\Psi_{\nu,\rho}(k).
  \label{eq:sc-determinant-factorization}
\end{equation}
The Wronskian also gives
\[
  \theta_\nu'(x)=\frac2{\pi xM_\nu(x)^2}>0.
\]
Nicholson's integral representation \cite[Eq.~(10.9.30)]{DLMF}
\[
  M_\nu(x)^2=\frac8{\pi^2}\int_0^\infty
  \cosh(2\nu t)K_0(2x\sinh t)\,dt
\]
and the strict decrease of \(K_0\) imply that \(M_\nu(x)^2\) is strictly
decreasing in \(x\).  Therefore
\begin{equation}
  \Psi_{\nu,\rho}'(k)
  =\frac2{\pi k}
  \left(M_\nu(k)^{-2}-M_\nu(\rho k)^{-2}\right)>0.
  \label{eq:sc-phase-monotonicity}
\end{equation}
The branch condition gives \(\Psi_{\nu,\rho}(0+)=0\), while the standard
large-argument Bessel expansion gives
\[
  \Psi_{\nu,\rho}(k)=(1-\rho)k+O(k^{-1})\longrightarrow\infty
\]
for fixed \(\nu\) and \(\rho\).  Hence the positive determinant roots are
uniquely indexed by
\[
  \Psi_{\ell+1/2,\rho}(k_{\ell,n})=n\pi,
  \qquad n=1,2,\ldots.
\]
Strict monotonicity gives
\[
  k_{\ell,n}<K
  \quad\Longleftrightarrow\quad
  n\pi<\Psi_{\ell+1/2,\rho}(K).
\]
The direct-sum decomposition supplies angular multiplicity \(2\ell+1\).
If the same numerical eigenvalue occurs at distinct values of \(\ell\), the
corresponding spherical-harmonic sectors are orthogonal and their
multiplicities add.  Finally,
\(\lambda_{\ell,n}\ge\ell(\ell+1)\), so only finitely many angular channels
can contribute below \(K^2\).  These observations prove
\eqref{eq:sc-exact-phase-count}, including every strict spectral wall.
\end{proof}

The only standard inputs in the preceding proof are completeness of the
spherical harmonics, regular one-dimensional Dirichlet Sturm--Liouville
theory, the Bessel equation, the Wronskian, Nicholson's representation, and
the usual small- and large-argument Bessel formulas; see \cite{DLMF}.

\subsection{Uniform phase enclosures and the strict lattice proxy}

For \(a>0\), define
\begin{equation}
  G_a(z)=
  \begin{cases}
  \displaystyle\frac1\pi\left(\sqrt{a^2-z^2}
  -z\arccos\frac za\right),&0\le z\le a,\\[1mm]
  0,&z>a.
  \end{cases}
  \label{eq:sc-model-action}
\end{equation}

\begin{lemma}[Elementary turning bound]
\label{lem:sc-turning-bound}
For \(a>0\) and \(0\leq z<a\),
\begin{equation}
  G_a(z)<\frac{(a-z)^{3/2}}{3\sqrt a}.
  \label{eq:sc-turning-point-bound}
\end{equation}
\end{lemma}

\begin{proof}
Write \(t=\cos\theta\).  Concavity of \(\sin\) on
\([0,\pi/4]\) gives
\(\sin(\theta/2)\geq\sqrt2\,\theta/\pi\), and hence
\[
  \arccos t\leq\frac\pi2\sqrt{1-t},
  \qquad 0\leq t\leq1.
\]
Since \(G_a(a)=0\) and
\[
  G_a(z)=\frac1\pi\int_z^a\arccos\frac ta\,dt,
\]
integration gives the stated strict inequality.
\end{proof}

For \(0\le z<a\), put
\[
  H_a(z)=\frac{3a^2+2z^2}{24\pi(a^2-z^2)^{3/2}},
\]
and define
\begin{equation}
  \mathcal F_a(z)=
  \begin{cases}
  \max\{G_a(z)-H_a(z),-1/4\},&0\le z<a,\\
  -1/4,&z\ge a.
  \end{cases}
  \label{eq:sc-auxiliary-F}
\end{equation}

We use the following published Bessel-phase result, stated here with every
hypothesis needed in the sequel.

\begin{lemma}[External uniform Bessel-phase enclosure]
\label{lem:sc-external-phase}
With the branch \eqref{eq:sc-phase-branch}, for every \(a>0\) and real
\(z\ge0\),
\begin{equation}
  \mathcal F_a(z)-\frac14
  <\frac{\theta_z(a)}\pi
  <G_a(z)-\frac14.
  \label{eq:sc-individual-phase-enclosure}
\end{equation}
If \(0<\mu<K\), \(0\le z\le K\), and
\[
  \Gamma_{K,\mu}(z)
  :=\frac{\theta_z(K)-\theta_z(\mu)}\pi,
\]
then
\begin{equation}
  0<\Gamma_{K,\mu}(z)
  <G_K(z)-\mathcal F_\mu(z)
  \le G_K(z)+\frac14.
  \label{eq:sc-global-difference-enclosure}
\end{equation}
If \(0\le z\le\mu\), including \(z=\mu\), then also
\begin{equation}
  G_K(z)-G_\mu(z)
  <\Gamma_{K,\mu}(z)
  <G_K(z)-G_\mu(z)+\frac14.
  \label{eq:sc-action-difference-enclosure}
\end{equation}
No value of the divergent expression \(H_\mu(\mu)\) is used in the last
statement.
\end{lemma}

Lemma~\ref{lem:sc-external-phase} is Theorem~5.1 and Lemma~5.2 of
\cite{FLPSannuli}, based on the real-order phase enclosures of
\cite{FLPSphase}.  The upper bound in
\eqref{eq:sc-global-difference-enclosure} follows immediately by subtracting
the lower bound for \(\theta_z(\mu)\) from the upper bound for
\(\theta_z(K)\) in \eqref{eq:sc-individual-phase-enclosure}; the correlated
two-sided estimate \eqref{eq:sc-action-difference-enclosure} is the sharper
part of the annulus result.  Both theorems are stated for real order, so the
half-integers \(z=\ell+\tfrac12\) require no interpolation.

Fix now \(K>0\), set \(\mu=\rho K\), and define
\begin{equation}
  A_{\rho,K}(z)=G_K(z)-G_\mu(z).
  \label{eq:sc-shell-action}
\end{equation}
For \(0\le z\le K\), let
\begin{equation}
  s_\mu(z)=
  \begin{cases}
  \min\{H_\mu(z),1/4\},&0\le z<\mu,\\
  1/4,&\mu\le z\le K,
  \end{cases}
  \qquad
  U_{\rho,K}(z)=A_{\rho,K}(z)+s_\mu(z).
  \label{eq:sc-optimized-proxy}
\end{equation}
The second line, rather than \(H_\mu\), is always used at \(z=\mu\).

\begin{proposition}[Strict phase-to-lattice reduction]
\label{prop:sc-phase-lattice}
For \(0<\rho<1\) and \(K>0\),
\begin{equation}
  \gamma_{\rho,K}(z)<U_{\rho,K}(z)
  \le A_{\rho,K}(z)+\frac14,
  \qquad 0\le z\le K.
  \label{eq:sc-pointwise-phase-proxy}
\end{equation}
Because every \(G_a\) is zero-extended, \(A_{\rho,K}\) is the piecewise
action \(G_K-G_{\rho K}\) below the inner interface and \(G_K\) at and
above it.  Consequently,
\begin{align}
  N_D(A_{\rho,1},K^2)
  &\le
  \sum_{\nu_\ell<K}2\nu_\ell
  [U_{\rho,K}(\nu_\ell)]_<
  \label{eq:sc-optimized-spectral-proxy}\\
  &\le
  S_{\rho,K}:=
  \sum_{\nu_\ell<K}2\nu_\ell
  \left\lfloor A_{\rho,K}(\nu_\ell)+\frac14\right\rfloor,
  \qquad \nu_\ell=\ell+\frac12.
  \label{eq:sc-coarse-spectral-proxy}
\end{align}
The cutoff is strict: an order \(\nu_\ell=K\) is inactive.  At a true phase
wall the radial count is one less than the ordinary floor; at a proxy wall
the ordinary floor in \eqref{eq:sc-coarse-spectral-proxy} is retained.
Moreover,
\begin{equation}
  \int_0^\infty2zA_{\rho,K}(z)\,dz
  =\frac{2}{9\pi}(1-\rho^3)K^3.
  \label{eq:sc-exact-action-integral}
\end{equation}
\end{proposition}

\begin{proof}
For \(0\le z<\mu\), the first upper bound in
\eqref{eq:sc-global-difference-enclosure} reads
\[
  \gamma_{\rho,K}(z)
  <A_{\rho,K}(z)+G_\mu(z)-\mathcal F_\mu(z).
\]
From \eqref{eq:sc-auxiliary-F},
\[
  G_\mu(z)-\mathcal F_\mu(z)
  =\min\{H_\mu(z),G_\mu(z)+1/4\}.
\]
The upper half of \eqref{eq:sc-action-difference-enclosure} also gives
\(\gamma_{\rho,K}(z)<A_{\rho,K}(z)+1/4\).  Taking the better of these two
strict upper bounds, and using \(G_\mu(z)+1/4\ge1/4\), gives
\[
  \gamma_{\rho,K}(z)
  <A_{\rho,K}(z)+\min\{H_\mu(z),1/4\}
  =U_{\rho,K}(z).
\]
For \(\mu\le z\le K\), one has \(G_\mu(z)=0\), and
\eqref{eq:sc-global-difference-enclosure} gives
\[
  \gamma_{\rho,K}(z)<G_K(z)+1/4=U_{\rho,K}(z).
\]

Orders \(z\ge K\) do not contribute.  Indeed,
\eqref{eq:sc-individual-phase-enclosure} gives
\[
  \frac{\theta_z(K)}\pi<-\frac14,
  \qquad
  \frac{\theta_z(\mu)}\pi>-\frac12,
\]
because \(G_K(z)=0\) and \(\mathcal F_\mu(z)=-1/4\).  Together with phase
monotonicity this yields
\[
  0<\gamma_{\rho,K}(z)<\frac14<1,
\]
so there is no positive phase level below \(K\).  This includes \(z=K\).

For positive \(x<y\), every positive integer strictly below \(x\) is also
strictly below \(y\), so \([x]_<\le[y]_<\le\lfloor y\rfloor\).  Applying
this observation to the exact phase count
\eqref{eq:sc-exact-phase-count}, and then using
\(s_\mu\le1/4\), proves
\eqref{eq:sc-optimized-spectral-proxy}--
\eqref{eq:sc-coarse-spectral-proxy} without changing either wall
convention.

Finally, scaling \(z=ax\) gives
\[
  \int_0^a2zG_a(z)\,dz
  =\frac{2a^3}{\pi}
  \left(\int_0^1x\sqrt{1-x^2}\,dx
  -\int_0^1x^2\arccos x\,dx\right).
\]
The two elementary integrals are \(1/3\) and \(2/9\), respectively; hence
the last expression is \(2a^3/(9\pi)\).  Subtracting the result with
\(a=\mu=\rho K\) from that with \(a=K\) proves
\eqref{eq:sc-exact-action-integral}.
\end{proof}

\section{Product and complementary-action tools}
\label{sec:product-action-tools}

We now prove the P\'olya inequality uniformly as the shell becomes thin.  In
this section the counting function is strict:
\[
  N_D(\Omega,K^2)=\#\{j:\lambda_j(\Omega)<K^2\}.
\]
We retain the global convention
\([u]_<:=\#\{n\in\mathbb N:n<u\}=\max\{0,\lceil u\rceil-1\}\).
Thus a spectral level lying exactly on the wall \(K^2\) is not counted.  Put
\begin{equation}
  \varepsilon=1-\rho,
  \qquad a=\varepsilon K,
  \qquad
  m_\varepsilon=1-\varepsilon+\frac{\varepsilon^2}{3}.
  \label{thin:eq-normalized-variables}
\end{equation}
Since \(1-\rho^3=3\varepsilon m_\varepsilon\), the desired right-hand side
has the two equivalent forms
\begin{equation}
  \frac{2}{9\pi}(1-\rho^3)K^3
  =\frac{1}{\varepsilon^2}\frac{2a^3}{3\pi}m_\varepsilon.
  \label{thin:eq-weyl-normalization}
\end{equation}

The proof uses the optimized phase estimate already established in
Proposition~\ref{prop:sc-phase-lattice}.  We record precisely the consequence
that is used here.  For
\(\nu_\ell=\ell+\tfrac12\), let
\(\gamma_{\rho,K}(\nu_\ell)\) be the normalized radial Bessel phase.  If
\(0\leq\nu_\ell\leq K\), that proposition gives
\begin{equation}
  \gamma_{\rho,K}(\nu_\ell)
  < A_{\rho,K}(\nu_\ell)+\frac14,
  \label{thin:eq-phase-input}
\end{equation}
where, after the scaling in \eqref{thin:eq-normalized-variables},
\(A_{\rho,K}(\nu)=\mathcal A(\varepsilon\nu)\) and
\(\mathcal A\) is the explicit action in
\eqref{thin:eq-action-definition} below.  The exact radial count in channel
\(\ell\) is
\begin{equation}
  \#\{n\geq1:n<\gamma_{\rho,K}(\nu_\ell)\}.
  \label{thin:eq-exact-phase-count}
\end{equation}
In particular, \eqref{thin:eq-phase-input} is an inequality for the strict
count; it is not a replacement of that count by an ordinary floor.

\subsection{The product comparison}

Under the standard unitary radial transformation, the angular channel
\(\ell\in\mathbb N_0\) is the Dirichlet operator
\[
  L_\ell=-\frac{d^2}{dr^2}+\frac{\ell(\ell+1)}{r^2}
  \quad\hbox{in }L^2((\rho,1),dr),
\]
and it occurs with multiplicity \(2\ell+1\).  The following lemma includes
the direction of the min--max comparison and every strict radial and angular
wall.

\begin{lemma}[Strict product majorant]
\label{thin:lem-product-majorant}
Let \(0<\varepsilon<1\), \(a=\varepsilon K\), and define
\begin{equation}
  N(a)=\max\left\{0,\left\lceil\frac a\pi\right\rceil-1\right\}.
  \label{thin:eq-radial-count}
\end{equation}
For \(1\leq n\leq N(a)\), put
\[
  b_n=\sqrt{a^2-n^2\pi^2},
  \qquad
  M_n=
  \left\lceil
  \sqrt{\left(\frac{b_n}{\varepsilon}\right)^2+\frac14}
  -\frac12
  \right\rceil.
\]
Then
\begin{equation}
  N_D(A_{\rho,1},K^2)
  \leq \mathcal P_\varepsilon(a)
  :=\sum_{n=1}^{N(a)}M_n^2.
  \label{thin:eq-product-majorant}
\end{equation}
In particular, the strict shell count is zero for \(0\leq a\leq\pi\).
\end{lemma}

\begin{proof}
The two quadratic forms have the common domain \(H_0^1(\rho,1)\), and
\[
  \int_\rho^1\left(
     |u'|^2+\frac{\ell(\ell+1)}{r^2}|u|^2
  \right)dr
  \geq
  \int_\rho^1\left(|u'|^2+\ell(\ell+1)|u|^2\right)dr,
\]
because \(r\leq1\).  The min--max principle therefore gives
\[
  \lambda_{\ell,n}
  \geq \left(\frac{n\pi}{\varepsilon}\right)^2
       +\ell(\ell+1),
  \qquad n\geq1.
\]
Consequently, a shell eigenvalue in this channel can be counted below
\(K^2\) only if
\begin{equation}
  n^2\pi^2+\varepsilon^2\ell(\ell+1)<a^2.
  \label{thin:eq-product-condition}
\end{equation}
The active radial integers are exactly those satisfying \(n\pi<a\), whose
number is \eqref{thin:eq-radial-count}.  For a fixed active \(n\), condition
\eqref{thin:eq-product-condition} is
\[
  \ell<
  \sqrt{\left(\frac{b_n}{\varepsilon}\right)^2+\frac14}-\frac12.
\]
Hence the allowed angular indices are
\(0,\ldots,M_n-1\), and their full spherical multiplicity is
\[
  \sum_{\ell=0}^{M_n-1}(2\ell+1)=M_n^2.
\]
Summing proves \eqref{thin:eq-product-majorant}.

If \(a=m\pi\), the radial level \(n=m\) is excluded by the strict
inequality \(n\pi<a\); thus \(N(a)=m-1\).  At an angular wall
\(b_n^2/\varepsilon^2=L(L+1)\), the displayed expression inside the ceiling
is exactly \(L\), so \(M_n=L\) and the level \(\ell=L\) is excluded.  The
jump of multiplicity \(2L+1\) occurs only above the wall.  Finally, no
positive integer satisfies \(n\pi<a\) when \(0\leq a\leq\pi\), including
both endpoints, so the product majorant and hence the shell count are zero
there.
\end{proof}


\subsection{The complementary-action comparison}

For \(r\in\{\rho,1\}\) and \(0\leq y\leq ar\), define
\begin{equation}
  J_r(y)=\sqrt{a^2r^2-y^2}
          -y\arccos\frac{y}{ar},
  \label{thin:eq-J-action}
\end{equation}
and set
\begin{equation}
  \mathcal A(y)=
  \begin{cases}
  \dfrac{J_1(y)-J_\rho(y)}{\pi\varepsilon},
    &0\leq y\leq\rho a,\\[6pt]
  \dfrac{J_1(y)}{\pi\varepsilon},
    &\rho a\leq y\leq a.
  \end{cases}
  \label{thin:eq-action-definition}
\end{equation}
We extend \(\mathcal A(y)=0\) for \(y>a\).

\begin{lemma}[Action geometry and curvature switch]
\label{thin:lem-action-geometry}
The function \(\mathcal A\) is continuous, continuously differentiable at
\(y=\rho a\), and strictly decreasing from
\[
  T:=\mathcal A(0)=\frac a\pi
  \quad\hbox{to}\quad \mathcal A(a)=0.
\]
Let \(R:[0,T]\to[0,a]\) be its decreasing inverse, let
\[
  \tau=\mathcal A(\rho a),
  \qquad F(t)=R(t)^2,
  \qquad g(t)=-F'(t).
\]
Then
\begin{equation}
\begin{gathered}
  g\text{ decreases on }(0,\tau),
  \qquad g\text{ increases on }(\tau,T),\\
  F(0)=a^2,
  \qquad F(\tau)=\rho^2a^2,
  \qquad F(T)=0,
  \qquad g(T)=2\pi\rho a,
\end{gathered}
  \label{thin:eq-action-geometry}
\end{equation}
where \(g(T)\) denotes the one-sided limit.  Moreover \(g(t)=O(t^{-1/3})\)
as \(t\downarrow0\).
\end{lemma}

\begin{proof}
Direct differentiation on the appropriate open intervals gives
\begin{equation}
  J_r'(y)=-\arccos\frac{y}{ar},
  \qquad
  J_r''(y)=\frac1{\sqrt{a^2r^2-y^2}}.
  \label{thin:eq-J-derivatives}
\end{equation}
At \(y=\rho a\), \(J_\rho(\rho a)=0\), and the two one-sided first
derivatives of \eqref{thin:eq-action-definition} both equal
\(-\arccos\rho/(\pi\varepsilon)\).  Thus the two pieces and their first
derivatives match.  Both derivatives are negative away from \(y=a\), so the
inverse exists on exactly the stated interval.

On the outer branch \(\rho a<y<a\), put
\(\alpha=\arccos(y/a)\).  If \(t=\mathcal A(y)\), then
\begin{equation}
  g(t)=-\frac{2y}{\mathcal A'(y)}
      =\frac{2\pi\varepsilon y}{\alpha}.
  \label{thin:eq-g-outer}
\end{equation}
The function \(y/\alpha\) increases with \(y\), whereas \(y=R(t)\)
decreases with \(t\); hence \(g\) decreases on \(0<t<\tau\).

On the inner branch \(0<y<\rho a\), put
\[
  \delta(y)=\arccos\frac ya-\arccos\frac{y}{\rho a}>0.
\]
Differentiating \(\arccos(y/(ar))\) with respect to \(r\) gives the exact
identity
\begin{equation}
  \frac{\delta(y)}{y}
  =\int_\rho^1\frac{dr}{r\sqrt{a^2r^2-y^2}}.
  \label{thin:eq-delta-integral}
\end{equation}
The integrand increases strictly with \(y\), so \(\delta(y)/y\) increases,
and \(y/\delta(y)\) decreases, with \(y\).  Since
\[
  g(t)=\frac{2\pi\varepsilon y}{\delta(y)}
\]
and \(y=R(t)\) decreases, \(g\) increases on \(\tau<t<T\).  Formula
\eqref{thin:eq-g-outer} and its inner counterpart agree at the interface.

The three values of \(F\) follow from
\(R(0)=a\), \(R(\tau)=\rho a\), and \(R(T)=0\).  As \(y\downarrow0\),
\eqref{thin:eq-delta-integral} gives
\[
  \delta(y)=\frac{\varepsilon}{\rho a}y+O(y^3),
\]
and therefore \(g(T)=2\pi\rho a\).

Finally, on the outer branch write \(y=a\cos\theta\).  Then
\begin{equation}
  t=\frac{a}{\pi\varepsilon}
       (\sin\theta-\theta\cos\theta)
    \sim\frac{a\theta^3}{3\pi\varepsilon},
  \qquad
  g(t)=\frac{2\pi\varepsilon a\cos\theta}{\theta}
       =O(t^{-1/3}).
  \label{thin:eq-improper-asymptotic}
\end{equation}
This proves the asserted improper endpoint behavior.  All square roots and
arccosines in the proof are taken on their stated closed domains; every
division occurs in an open interval where its denominator is positive, and
the endpoint zeros are handled by the limits just computed.
\end{proof}

\begin{lemma}[Strict phase-to-action layer cake]
\label{thin:lem-layer-cake}
For \(\ell\in\mathbb N_0\), set
\begin{equation}
  y_\ell=\varepsilon\left(\ell+\frac12\right),
  \qquad
  q_\ell=\left[\mathcal A(y_\ell)+\frac14\right]_<,
  \qquad
  P_{\mathcal A}=\varepsilon\sum_{\ell\geq0}y_\ell q_\ell.
  \label{thin:eq-phase-proxy}
\end{equation}
Then
\begin{equation}
  N_D(A_{\rho,1},K^2)
  \leq\frac{2P_{\mathcal A}}{\varepsilon^2}.
  \label{thin:eq-phase-transfer}
\end{equation}
For \(n\geq1\), put \(x_n=n-\tfrac14\), and define
\begin{equation}
  M_\varepsilon(x)
  =\#\left\{\ell\in\mathbb N_0:
       \varepsilon\left(\ell+\frac12\right)<x\right\}
  =\max\left\{0,
       \left\lceil\frac{x}{\varepsilon}-\frac12\right\rceil
       \right\}.
  \label{thin:eq-half-integer-count}
\end{equation}
The proxy has the exact finite representation
\begin{equation}
  P_{\mathcal A}
  =\frac{\varepsilon^2}{2}
    \sum_{\substack{n\geq1\\x_n<T}}
       M_\varepsilon(R(x_n))^2.
  \label{thin:eq-discrete-layer-cake}
\end{equation}
Moreover,
\begin{equation}
  I:=\frac12\int_0^T F(t)\,dt
  =\frac{a^3}{3\pi}
     \left(1-\varepsilon+\frac{\varepsilon^2}{3}\right),
  \qquad
  \frac{2I}{\varepsilon^2}
  =\frac{2}{9\pi}(1-\rho^3)K^3.
  \label{thin:eq-continuous-normalization}
\end{equation}
\end{lemma}

\begin{proof}
For \(y_\ell<a\), equations \eqref{thin:eq-phase-input} and
\eqref{thin:eq-exact-phase-count} show that the number of radial levels in
channel \(\ell\) is at most \(q_\ell\).  If \(y_\ell\geq a\), equivalently
\(\nu_\ell\geq K\), the product comparison gives
\[
  \lambda_{\ell,1}
  \geq\frac{\pi^2}{\varepsilon^2}+\ell(\ell+1)
  =\frac{\pi^2}{\varepsilon^2}+\nu_\ell^2-\frac14
  >K^2,
\]
so there is no radial level; at \(y_\ell=a\), the definition also gives
\(q_\ell=[1/4]_< =0\).  Multiplicity \(2\ell+1=2y_\ell/\varepsilon\)
now yields
\[
  N_D(A_{\rho,1},K^2)
  \leq\sum_{\ell\geq0}\frac{2y_\ell}{\varepsilon}q_\ell
  =\frac{2P_{\mathcal A}}{\varepsilon^2},
\]
which proves \eqref{thin:eq-phase-transfer}.  If
\(\mathcal A(y_\ell)+1/4=m\in\mathbb Z\), then \(q_\ell=m-1\), so the
phase-proxy wall is excluded exactly as required.

For every \(y\geq0\), the strict bracket is
\[
  \left[\mathcal A(y)+\frac14\right]_<
  =\#\{n\geq1:x_n<\mathcal A(y)\}.
\]
The sums are finite.  Since \(\mathcal A\) and \(R\) are decreasing,
\(x_n<\mathcal A(y_\ell)\) is equivalent to \(y_\ell<R(x_n)\).  Therefore
\[
\begin{aligned}
  P_{\mathcal A}
  &=\varepsilon
    \sum_{\substack{n\geq1\\x_n<T}}
      \sum_{y_\ell<R(x_n)}y_\ell \\
  &=\varepsilon
    \sum_{\substack{n\geq1\\x_n<T}}
      \varepsilon\sum_{\ell=0}^{M_\varepsilon(R(x_n))-1}
        \left(\ell+\frac12\right) \\
  &=\frac{\varepsilon^2}{2}
    \sum_{\substack{n\geq1\\x_n<T}}
      M_\varepsilon(R(x_n))^2,
\end{aligned}
\]
because \(\sum_{\ell=0}^{M-1}(\ell+1/2)=M^2/2\).  This proves
\eqref{thin:eq-discrete-layer-cake}.  Formula
\eqref{thin:eq-half-integer-count} follows from the same strict counting;
at \(x/\varepsilon=m+1/2\), it equals \(m\), excluding \(\ell=m\).

Differentiation of \eqref{thin:eq-J-action} with respect to \(r\) gives
\[
  \frac{\partial J_r(y)}{\partial r}
  =\sqrt{a^2-\frac{y^2}{r^2}}
  \quad(0\leq y\leq ar).
\]
Consequently the two cases of \eqref{thin:eq-action-definition} combine as
\begin{equation}
  \mathcal A(y)
  =\frac1{\pi\varepsilon}
    \int_\rho^1
      \left(a^2-\frac{y^2}{r^2}\right)_+^{1/2}\,dr,
  \label{thin:eq-action-integral}
\end{equation}
where \(s_+=\max\{s,0\}\).
The area identity for a decreasing inverse and Tonelli's theorem give
\[
  \frac12\int_0^T R(t)^2\,dt
  =\int_0^a y\mathcal A(y)\,dy.
\]
Using \eqref{thin:eq-action-integral} and then \(y=rz\),
\[
\begin{aligned}
  \int_0^a y\mathcal A(y)\,dy
  &=\frac1{\pi\varepsilon}
    \int_\rho^1r^2\,dr
    \int_0^a z\sqrt{a^2-z^2}\,dz \\
  &=\frac{a^3}{3\pi\varepsilon}\int_\rho^1r^2\,dr
   =\frac{a^3}{3\pi}
      \left(1-\varepsilon+\frac{\varepsilon^2}{3}\right).
\end{aligned}
\]
This is \eqref{thin:eq-continuous-normalization}; its final equality follows
from \(1-\rho^3=3\varepsilon m_\varepsilon\) and \(K=a/\varepsilon\).
\end{proof}

\begin{lemma}[Shifted sawtooth and radial quadrature deficit]
\label{thin:lem-radial-deficit}
Let
\[
  C(t)=\#\{n\geq1:n-\tfrac14<t\},
  \qquad w(t)=t-C(t),
  \qquad \mathfrak W(t)=\int_0^t w(s)\,ds,
\]
and put \(\Psi(t)=\mathfrak W(t)-t/4\).  Then, for all \(t\geq0\),
\begin{equation}
  \mathfrak W(t)\geq0,
  \qquad
  -\frac1{32}\leq\Psi(t)\leq\frac3{32}.
  \label{thin:eq-sawtooth-bounds}
\end{equation}
Furthermore, with
\begin{equation}
  D_{\mathrm{rad}}
  =\int_0^T F(t)\,dt
   -\sum_{\substack{n\geq1\\x_n<T}}F(x_n),
  \label{thin:eq-radial-deficit-definition}
\end{equation}
one has
\begin{equation}
  D_{\mathrm{rad}}
  \geq\frac{\rho^2a^2}{4}-\frac{\pi\rho a}{4}.
  \label{thin:eq-radial-deficit-bound}
\end{equation}
\end{lemma}

\begin{proof}
On \(0\leq t\leq3/4\), \(C(t)=0\), and hence
\begin{equation}
  \Psi(t)=\frac12\left(t-\frac14\right)^2-\frac1{32}.
  \label{thin:eq-first-sawtooth-cell}
\end{equation}
For \(x_n=n-1/4\), \(n\geq1\), and \(t=x_n+u\) with
\(0\leq u\leq1\), direct integration gives
\begin{equation}
  \Psi(t)=\frac12\left(u-\frac12\right)^2-\frac1{32}.
  \label{thin:eq-periodic-sawtooth-cell}
\end{equation}
These formulas agree at cell endpoints and prove the two-sided bound for
\(\Psi\).  On the first cell \(\mathfrak W=t^2/2\geq0\).  Thereafter
\(\mathfrak W=t/4+\Psi\geq t/4-1/32>0\), proving
\eqref{thin:eq-sawtooth-bounds}.  At a jump \(x_n\), strict counting gives
the one-sided values \(w(x_n^-)=3/4\) and \(w(x_n^+)=-1/4\), while
\(\mathfrak W\)
and \(\Psi\) are continuous.

Distributionally, \(dw=dt-dC\).  Since \(F(T)=0\), an atom at \(x_n=T\)
has zero weight and is consistent with its exclusion from
\eqref{thin:eq-radial-deficit-definition}.  Stieltjes integration by parts,
first on \([\eta,T]\) and then with \(\eta\downarrow0\), gives
\begin{equation}
  D_{\mathrm{rad}}=\int_0^T w(t)g(t)\,dt.
  \label{thin:eq-stieltjes-start}
\end{equation}
Indeed, \eqref{thin:eq-improper-asymptotic}, together with
\(w(t)=t\) and \(\mathfrak W(t)=t^2/2\) near zero, gives
\(w(t)g(t)=O(t^{2/3})\) and
\(\mathfrak W(t)g(t)=O(t^{5/3})\to0\), so the improper
boundary term vanishes.

Split \eqref{thin:eq-stieltjes-start} at the actual value \(\tau\), with
no assumption that \(\tau\) belongs to any particular sawtooth cell.  On
\((0,\tau)\), \(g\) decreases and \(\mathfrak W\geq0\); hence
\begin{equation}
  \int_0^\tau wg
  =\mathfrak W(\tau)g(\tau)-\int_0^\tau \mathfrak W\,dg
  \geq \mathfrak W(\tau)g(\tau).
  \label{thin:eq-decreasing-stieltjes}
\end{equation}
On \((\tau,T)\), \(g\) increases and \(w=1/4+\Psi'\) almost everywhere.
Adding its integration-by-parts identity to
\eqref{thin:eq-decreasing-stieltjes}, and using
\(\mathfrak W(\tau)=\tau/4+\Psi(\tau)\), cancels the interface value and
yields
\[
  D_{\mathrm{rad}}
  \geq \frac{F(\tau)}4+\frac{\tau g(\tau)}4
     +\Psi(T)g(T)-\int_\tau^T\Psi\,dg.
\]
Because \(dg\) is a positive measure on \((\tau,T)\), the bounds
\eqref{thin:eq-sawtooth-bounds} imply
\[
\begin{aligned}
  D_{\mathrm{rad}}
  &\geq \frac{F(\tau)}4-\frac{g(T)}8
    +g(\tau)\left(\frac\tau4+\frac3{32}\right) \\
  &\geq \frac{F(\tau)}4-\frac{g(T)}8.
\end{aligned}
\]
The two discarded quantities,
\(-\int_0^\tau \mathfrak W\,dg\) and
\(g(\tau)(\tau/4+3/32)\), are explicitly nonnegative.  Substitution of
\(F(\tau)=\rho^2a^2\) and \(g(T)=2\pi\rho a\) from
Lemma~\ref{thin:lem-action-geometry} proves
\eqref{thin:eq-radial-deficit-bound}.  Continuity of
\(\mathfrak W,\Psi,g\) at the
split shows that the same proof covers \(\tau\) below, on, or above every
sawtooth grid point.
\end{proof}


\section{Global two-ratio closure}
\label{sec:global-two-owner}

\subsection{Ratio-sharp defect theorem below the optical strip}

For $a>0$, let
\[
 G_a(z):=\frac1\pi\left(\sqrt{a^2-z^2}
 -z\arccos\frac za\right)\quad(0\leq z<a),
 \qquad G_a(z):=0\quad(z\geq a),
\]
and put
\[
 A(z):=G_K(z)-G_{\rho K}(z),
 \qquad T:=\frac{(1-\rho)K}{\pi}.
\]
The strict phase reduction and the inverse-action layer cake give
\begin{equation}
 N_D(A_{\rho,1},K^2)\leq P,\qquad
 P=\sum_{n-1/4<T}M(R(n-1/4))^2,\qquad
 W=\int_0^TR(t)^2\,dt,
 \label{rs-summary:eq:layer-cake}
\end{equation}
where $R$ is the decreasing inverse of $A$ and $M(r)$ counts
half-integers strictly below $r$.  Thus
\[
 W-P=\mathcal D_{\rm rad}-\mathcal E_{\rm ang}.
\]

Write $\theta=\arccos\rho$ and
$N=\#\{n\geq1:n-1/4<T\}$.  The exact action derivative satisfies
$-A'\leq\theta/\pi$.  One disjoint left block per active layer therefore
gives the ratio-sharp estimate
\begin{equation}
 \mathcal E_{\rm ang}
 <\frac{(1-\rho^2)K^2}{6}
 -N\left(\frac{3\pi}{8\theta}-\frac14\right).
 \label{rs-summary:eq:angular}
\end{equation}
Strictness includes a half-integer angular wall.

For $0<\rho\leq39/50$ and $K\geq\pi/(1-\rho)$, the curvature interface
$\tau=KG_1(\rho)$ satisfies $\tau>1/4$.  Replacing the shifted-sawtooth
primitive by its $C^1$ tangent envelope
\[
 L_\sharp(t)=
 \begin{cases}t^2/2,&t\leq1/4,\\t/4-1/32,&t\geq1/4,\end{cases}
\]
and retaining the decreasing-branch Stieltjes remainder gives
\begin{equation}
 \mathcal D_{\rm rad}\geq
 B_0(K)-\frac{\pi\rho K}{4(1-\rho)}
 +\frac{\pi\rho K}{4\arccos\rho},
 \qquad
 B_0(K):=\int_0^{1/4}R_B(t)^2\,dt,
 \label{rs-summary:eq:radial}
\end{equation}
where $R_B$ is the inverse ball action.  The turning estimate
$G_1(s)\geq2\sqrt2(1-s)^{3/2}/(3\pi)$ yields
\begin{equation}
 B_0(K)\geq\frac{K^2}{4}-\alpha K^{4/3}+\beta K^{2/3},
 \quad
 \alpha=\frac{6C}{5\,4^{5/3}},\quad
 \beta=\frac{3C^2}{7\,4^{7/3}},\quad
 C=\frac{(3\pi)^{2/3}}2.
 \label{rs-summary:eq:ball}
\end{equation}

Since $N\geq1$ and
$\arccos\rho\leq(\pi/2)\sqrt{1-\rho}$, it remains to prove positivity of
\begin{equation}
 \begin{aligned}
 \underline{\mathcal G}(\rho,K):={}&
 \frac{1+2\rho^2}{12}K^2-\alpha K^{4/3}+\beta K^{2/3}
 -\frac{\pi\rho K}{4(1-\rho)}\\
 &+\frac{\rho K}{2\sqrt{1-\rho}}
 +\frac3{4\sqrt{1-\rho}}-\frac14.
 \end{aligned}
 \label{rs-summary:eq:scalar}
\end{equation}
On the base $K=\pi/(1-\rho)$, the substitution
$x=(1-\rho)^{-1/6}$ reduces \eqref{rs-summary:eq:scalar} to one explicit
degree-nine expression.  The detailed support module proves it positive
from $157/50<\pi<22/7$, positive-side cubing, and
\[
 F(1)>\frac14,\qquad F'(1)>-\frac52,\qquad F''(x)>14.
\]
It also proves that the base $K$-derivative is positive and
\[
 \partial_K^2\underline{\mathcal G}
 >\frac{47447}{443682}>0.
\]
These are finite exact inequalities, not sampled or interval premises.

\begin{theorem}[Global low- and middle-ratio theorem]
\label{rs:thm:global-proxy}
For
\[
 0<\rho\leq\frac{39}{50},
 \qquad K\geq\frac\pi{1-\rho},
\]
one has
\[
 N_D(A_{\rho,1},K^2)\leq P<W(\rho,K).
\]
\end{theorem}

\begin{proof}
Equations \eqref{rs-summary:eq:angular}--\eqref{rs-summary:eq:ball}
give
\[
 W-P>\underline{\mathcal G}(\rho,K).
\]
The exact scalar closure just described proves the right side positive on
the stated domain.  All intermediate identities and the rational
convexity proof are given in the accompanying analytic supplement.
\end{proof}


\subsection{All-frequency optical theorem}

\begin{lemma}[All-frequency optical theorem]
\label{cm:lem-optical}
For \(39/50\leq\rho<1\) and \(K\geq0\),
\(N_D(A_{\rho,1},K^2)\leq W(\rho,K)\), with strict inequality for
\(K>0\).
\end{lemma}

\begin{proof}
Put \(\varepsilon=1-\rho\leq11/50\),
\(a=\varepsilon K\), and \(c=1126/625\).  We use two inclusive screens
meeting at \(a=c/\varepsilon\).

For the product screen, the strict min--max condition
\eqref{thin:eq-product-condition} implies that a counted mode satisfies
\[
 n^2\pi^2+\varepsilon^2\ell(\ell+1)<a^2.
\]
There is no mode for \(a\leq\pi\).  For \(a>\pi\), write
\(a/\pi=N+t\), with \(N\geq1\), \(0<t\leq1\), assigning
\((N,t)=(m-1,1)\) at \(a=m\pi\).  If
\[
 b_n=\sqrt{a^2-n^2\pi^2},\qquad
 M_n=\left\lceil\sqrt{(b_n/\varepsilon)^2+1/4}-1/2\right\rceil,
\]
the exact angular multiplicity is bounded by
\[
 \mathcal P_\varepsilon(a)=\sum_{n=1}^N M_n^2.
\]
Set
\[
 I_0(a)=\int_0^{a/\pi}(a^2-\pi^2x^2)\,dx
       =\frac{2a^3}{3\pi},
 \qquad
 S_0(a)=\sum_{n=1}^N b_n^2,
 \qquad D(a)=I_0(a)-S_0(a).
\]
Direct summation of the continuous-minus-discrete product deficit gives
\begin{equation}
 \frac{D(a)}{\pi^2}=\frac{N^2}{2}
 +N\left(t^2+\frac16\right)+\frac{2t^3}{3}.
 \label{cm:eq-optical-product-deficit}
\end{equation}
For \(C_D=1382/3125\), the difference
\(D/\pi^2-C_D(N+t)^2\) increases in \(N\), and at \(N=1\) is minimized
at \(t=C_D\), where it equals
\(1822532/91552734375>0\).  Thus \(D>C_Da^2\).

For \(x>0\),
\[
\left\lceil\sqrt{x^2+\frac14}-\frac12\right\rceil^2
 <x^2+x+1.
\]
Indeed the ceiling is strictly smaller than
\(\sqrt{x^2+1/4}+1/2\), whose square is
\(x^2+1/2+\sqrt{x^2+1/4}<x^2+x+1\).  Applying this with
\(x=b_n/\varepsilon\), then using the monotone left-sum estimate and the
quarter-circle integral,
\[
 \sum_{n=1}^N b_n
 <\int_0^{a/\pi}\sqrt{a^2-\pi^2x^2}\,dx=\frac{a^2}{4},
 \qquad N<\frac a\pi,
\]
gives
\[
 \varepsilon^2\mathcal P_\varepsilon
 <S_0+\varepsilon a^2/4+\varepsilon^2a/\pi.
\]
On \(a\leq c/\varepsilon\), using \(q=106/333>1/\pi\) and
\(\varepsilon\leq11/50\), the exact reserve is
\begin{equation}
 C_D-\left(\frac{2cq}{3}+\frac{11}{200}
 +(11/50)^2q^2\right)
 =\frac{39569}{2772225000}>0.
 \label{cm:eq-optical-low-reserve}
\end{equation}
Indeed, \(a\leq c/\varepsilon\), \(\varepsilon\leq11/50\),
\(a>\pi\), and \(q>1/\pi\) give
\[
 \varepsilon I_0+\frac{\varepsilon a^2}{4}
 +\frac{\varepsilon^2a}{\pi}
 <a^2\left(\frac{2cq}{3}+\frac{11}{200}
 +(11/50)^2q^2\right)<C_Da^2<D.
\]
Consequently
\[
 \varepsilon^2\mathcal P_\varepsilon
 <I_0-D+\frac{\varepsilon a^2}{4}
 +\frac{\varepsilon^2a}{\pi}
 <(1-\varepsilon)I_0<m_\varepsilon I_0
 =\varepsilon^2W,
\]
where the last identity is \eqref{thin:eq-weyl-normalization}.
This proves the low screen, including all radial and angular equality walls.

For the complementary-action screen, use the inverse-action construction
proved in Lemma~\ref{thin:lem-action-geometry}: let the decreasing shell
action have inverse \(R_{\mathcal A}\), put \(F=R_{\mathcal A}^2\), and
\(g=-F'\).  If \(\tau\) is the actual curvature interface and
\(T=a/\pi\), that lemma gives
\[
 g\text{ decreasing on }(0,\tau),\quad
 g\text{ increasing on }(\tau,T),
\]
\[
 F(0)=a^2,\quad F(\tau)=\rho^2a^2,\quad F(T)=0,\quad
 g(T)=2\pi\rho a.
\]
For \(C(t)=\#\{n:n-1/4<t\}\), set \(w=t-C(t)\) and
\(\Psi(t)=\int_0^tw(s)\,ds-t/4\).  The complete sawtooth and Stieltjes
calculation in Lemma~\ref{thin:lem-radial-deficit}, including the ungridded
curvature interface \(\tau\), gives
\begin{equation}
 D_{\rm rad}:=\int_0^TF(t)\,dt
 -\sum_{n-1/4<T}F(n-1/4)
 \geq\frac{\rho^2a^2}{4}-\frac{\pi\rho a}{4}.
 \label{cm:eq-optical-radial-deficit}
\end{equation}
The exact half-integer angular count from
Lemma~\ref{thin:lem-layer-cake} satisfies
\(\varepsilon^2M_\varepsilon(x)^2<(x+\varepsilon/2)^2\), even on an
angular wall.  Since fewer than \(a/\pi+1/4\) radial layers occur, the
angular error is strictly below
\begin{equation}
 E=\left(\frac a\pi+\frac14\right)
 \left(\varepsilon a+\frac{\varepsilon^2}{4}\right).
 \label{cm:eq-optical-angular-error}
\end{equation}
On \(a\geq c/\varepsilon\), one has \(1/a\leq\varepsilon/c\).  Put
\[
 L_{\rm opt}(\varepsilon)
 :=\frac{(1-\varepsilon)^2}{4}
 -\frac{(22/7)(1-\varepsilon)\varepsilon}{4c},
\]
\[
 A_{\rm opt}(\varepsilon)
 :=q\varepsilon+\frac{\varepsilon^2}{4c}
 +\frac{q\varepsilon^3}{4c}
 +\frac{\varepsilon^4}{16c^2}.
\]
Dividing \eqref{cm:eq-optical-radial-deficit} and
\eqref{cm:eq-optical-angular-error} by \(a^2\), then using
\(\pi<22/7\) and \(q>1/\pi\), gives
\[
 \frac{D_{\rm rad}}{a^2}\geq L_{\rm opt}(\varepsilon),
 \qquad
 \frac{E}{a^2}<A_{\rm opt}(\varepsilon).
\]
For \(0<\varepsilon\leq11/50<1/2\),
\[
 L_{\rm opt}'(\varepsilon)
 =-\frac{1-\varepsilon}{2}
 -\frac{22/7}{4c}(1-2\varepsilon)<0,
 \qquad
 A_{\rm opt}'(\varepsilon)>0.
\]
At \(\varepsilon=11/50\), the exact endpoint difference is
\begin{equation}
 L_{\rm opt}(11/50)-A_{\rm opt}(11/50)
 =\frac{14817541}{472867032960000}>0.
 \label{cm:eq-optical-high-reserve}
\end{equation}
Thus the strict layer-cake sum is below the action integral, and the phase
bridge proves the high screen.  The common face lies in the nonzero radial
range because
\[
 \frac c\varepsilon\geq\frac{c}{11/50}
 =\frac{2252}{275}>\frac{22}{7}>\pi,
 \qquad
 \frac{2252}{275}-\frac{22}{7}=\frac{9714}{1925}>0.
\]
Both screens include \(a=c/\varepsilon\), so their union has neither a gap
nor an unowned equality face.  At \(K=0\) both sides vanish.
\end{proof}

\section{Proof of the spectral theorem}

We now assemble the three disjoint phase-space owners.  Fix
\(0<\rho<1\) and put
\[
 K_0(\rho):=\frac{\pi}{1-\rho}.
\]
If \(0\leq K\leq K_0(\rho)\), the strict product majorant in
Lemma~\ref{thin:lem-product-majorant} gives
\begin{equation}
 N_D(A_{\rho,1},K^2)=0.
 \label{eq:global-no-mode}
\end{equation}
This includes \(K=K_0(\rho)\): the first one-dimensional Dirichlet
frequency lies on the spectral wall and is excluded.

It remains to take \(K>K_0(\rho)\).  If
\(0<\rho<39/50\), apply
Theorem~\ref{rs:thm:global-proxy}.  If
\(39/50\leq\rho<1\), apply Lemma~\ref{cm:lem-optical}.  These cases,
together with \eqref{eq:global-no-mode}, are disjoint and exhaustive.
Therefore
\begin{equation}
  N_D(A_{\rho,1},K^2)
  \le\frac{2}{9\pi}(1-\rho^3)K^3
  \qquad(K\ge0).
  \label{eq:unit-shell-theorem}
\end{equation}

Now
\[
  \omega_3=\frac{4\pi}{3},
  \qquad
  L_3=\frac{\omega_3}{(2\pi)^3}=\frac1{6\pi^2},
  \qquad
  |A_{\rho,1}|=\frac{4\pi}{3}(1-\rho^3),
\]
so
\begin{equation}
  L_3|A_{\rho,1}|=\frac{2}{9\pi}(1-\rho^3).
  \label{eq:weyl-normalization}
\end{equation}
Taking \(K=\sqrt\Lambda\) in \eqref{eq:unit-shell-theorem} proves the
unit-shell form of Theorem~\ref{thm:spectral} for every \(\Lambda\ge0\).

For arbitrary \(0<r<R\), put \(\rho=r/R\).  The unitary dilation
\[
  (D_Ru)(x)=R^{-3/2}u(x/R)
\]
maps \(L^2(A_{\rho,1})\) onto \(L^2(A_{r,R})\), and its Dirichlet form
scales by \(R^{-2}\).  Hence
\[
  \lambda_j^D(A_{r,R})=R^{-2}\lambda_j^D(A_{\rho,1}),
\]
and the strict threshold transforms as
\begin{equation}
  N_D(A_{r,R},\Lambda)
  =N_D(A_{\rho,1},R^2\Lambda).
  \label{eq:dilation-count}
\end{equation}
Applying the unit-shell result at \(R^2\Lambda\),
\[
\begin{aligned}
  N_D(A_{r,R},\Lambda)
  &\le\frac{2}{9\pi}(1-\rho^3)(R^2\Lambda)^{3/2}\\
  &=\frac{2}{9\pi}(R^3-r^3)\Lambda^{3/2}\\
  &=L_3|A_{r,R}|\Lambda^{3/2}.
\end{aligned}
\]
This proves Theorem~\ref{thm:spectral} for every \(R>0\), including
\(R<1\).

\section{Proof that spherical shells do not tile}

We prove Theorem~\ref{thm:nontiling} independently of the spectral argument.
Set \(\alpha=r/R\in(0,1)\) and scale by \(1/R\).  It suffices to consider
\[
  T=A_{\alpha,1}=\{x:\alpha<|x|<1\}.
\]
Every rigid motion has the form \(x\mapsto Qx+a\), with \(Q\) orthogonal.
Because \(T\) is radial, \(Q(T)=T\), so every congruent copy is a translate
\[
  T_a=a+T=\{x:\alpha<|x-a|<1\}.
\]

Assume for contradiction that \(\{T_a:a\in\mathfrak C\}\) has
pairwise-disjoint interiors and covers almost everywhere.  Duplicate centers
are impossible because they give the same nonempty open tile.

Fix a unit vector \(e\) and put
\[
  c=\frac{1+\alpha}{2}e,
  \qquad
  \delta=\frac{1-\alpha}{2}.
\]
Write \(B(x,s)=\{y\in\R^3:|y-x|<s\}\) for the open Euclidean ball.
Then \(B(c,\delta)\subset T\).  Therefore \(T_a\) contains
\(B(a+c,\delta)\), and disjointness implies
\begin{equation}
  |a-b|\ge1-\alpha
  \qquad(a\ne b).
  \label{eq:center-separation}
\end{equation}
Thus the centers are locally finite: a bounded set cannot contain infinitely
many points with fixed positive separation.  They are also countable, since
\[
  \mathfrak C
  =\bigcup_{m\ge1}(\mathfrak C\cap\overline B(0,m))
\]
is a countable union of finite sets.

Each tile boundary is a union of two spheres and is three-dimensional null.
The countable union
\[
  Z=\bigcup_{a\in\mathfrak C}\partial T_a
\]
is therefore null.  Exact open-shell coverage is already a special case of
almost-everywhere coverage.  If closed shells cover exactly or almost
everywhere, removing \(Z\) shows that their open interiors still cover almost
everywhere.  Hence it is enough to assume
\begin{equation}
  \left|\R^3\setminus\bigcup_{a\in\mathfrak C}T_a\right|=0.
  \label{eq:ae-coverage}
\end{equation}

Choose one tile \(T_{a_0}\) and its outer sphere
\[
  \Sigma=\{x:|x-a_0|=1\}.
\]
Fix \(p\in\Sigma\) and put \(u=p-a_0\).  For \(n\ge1\), let
\[
  U_n=B\left(p+\frac2n u,\frac1n\right).
\]
Every \(U_n\) lies outside \(T_{a_0}\), is open, and has positive volume.
By \eqref{eq:ae-coverage}, choose \(x_n\in U_n\cap T_{a_n}\) with
\(a_n\ne a_0\).  Since \(x_n\to p\) and \(|x_n-a_n|<1\), the centers
\(a_n\) eventually lie in a fixed bounded set.  Local finiteness lets us
pass to a subsequence with one fixed center \(a\ne a_0\).  Thus
\(p\in\overline{T_a}\).

The point \(p\) is not in \(T_a\).  Otherwise some
\(B(p,\varepsilon)\) lies in \(T_a\).  Choose
\[
  0<t<\min\{\varepsilon,1-\alpha\}.
\]
Then \(p-tu\in T_{a_0}\cap T_a\), contradicting disjoint interiors.
Therefore every \(p\in\Sigma\) belongs to the boundary of a neighboring
tile.

If \(p\in\Sigma\cap\partial T_a\), then
\(|p-a|\in\{\alpha,1\}\), hence \(|a-a_0|\le2\).  Only finitely many
centers satisfy this bound.  Consequently \(\Sigma\) is covered by finitely
many neighboring inner or outer boundary spheres.

Translate so that \(a_0=0\).  The intersection of \(\Sigma\) with a
neighboring sphere \(\{x:|x-a|=s\}\), \(s\in\{\alpha,1\}\), satisfies
\[
  |x|^2=1,
  \qquad |x-a|^2=s^2,
\]
and hence
\begin{equation}
  2a\cdot x=|a|^2+1-s^2.
  \label{eq:sphere-plane}
\end{equation}
Because \(a\ne0\), \eqref{eq:sphere-plane} is a genuine plane.  Its
intersection with \(\Sigma\) is empty, a tangent point, or a circle, and
therefore has zero surface measure on \(\Sigma\).  A finite union of such
sets cannot cover \(\Sigma\), whose area is \(4\pi\).

There is no coincident-sphere exception.  A neighboring outer sphere equals
\(\Sigma\) only when it has the same center, which is the original tile or a
forbidden duplicate.  A neighboring inner sphere cannot equal \(\Sigma\)
because \(\alpha<1\).  We have reached a contradiction.  Scaling back proves
Theorem~\ref{thm:nontiling}.

\appendix

\section{External inputs and verification boundary}
\label{app:external-boundary}

The proof above is self-contained with respect to every shell-specific
spectral reduction, weighted count, endpoint estimate, regional subtraction,
fixed-channel variational comparison, boundary assignment, and finite
exact-rational predicate.  It uses the following conventional external results, each only
in the scope stated where it enters the proof:
\begin{enumerate}
  \item standard spherical-harmonic decomposition, regular
        Sturm--Liouville theory, the contiguous-order zero interlacing
        in DLMF equation~(10.21.2), and the classical Bessel identities
        collected in \cite{DLMF};
  \item the uniform one-variable Bessel-phase enclosure of
        \cite{FLPSphase}, together with the explicitly quoted floor-sum
        inequalities of \cite{FLPSballs,FLPSannuli};
  \item inequalities~(i) and (ii), the two strict lower bounds in Lorch's
        published abstract, for the first positive zero of \(J_\nu\) on
        \(-1<\nu<\infty\), in precisely the half-integer specializations
        derived in the text \cite{Lorch1993}.
\end{enumerate}

All shell-specific finite decisions are exposed in the accompanying
\emph{Purely Analytic Supplement}.  They require only exact integer and
rational arithmetic, elementary integral inequalities, and the classical
real functions already present in the analytic argument.  The archived
interval programs have no place in this verification boundary; they are
optional regression checks only.

\begin{thebibliography}{5}

\bibitem{FLPSballs}
N.~Filonov, M.~Levitin, I.~Polterovich, and D.~A.~Sher,
\newblock P\'olya's conjecture for Euclidean balls,
\newblock \emph{Invent. Math.} \textbf{234} (2023), 129--169.
\newblock \href{https://doi.org/10.1007/s00222-023-01198-1}{doi:10.1007/s00222-023-01198-1}.

\bibitem{FLPSphase}
N.~Filonov, M.~Levitin, I.~Polterovich, and D.~A.~Sher,
\newblock Uniform enclosures for the phase and zeros of Bessel functions and
their derivatives,
\newblock \emph{SIAM J. Math. Anal.} \textbf{56} (2024), no.~6, 7644--7682.
\newblock \href{https://doi.org/10.1137/24M1642032}{doi:10.1137/24M1642032}.

\bibitem{FLPSannuli}
N.~Filonov, M.~Levitin, I.~Polterovich, and D.~A.~Sher,
\newblock P\'olya's conjecture for Dirichlet eigenvalues of annuli,
\newblock \emph{J. Lond. Math. Soc.} \textbf{113} (2026), no.~2, e70425.
\newblock \href{https://doi.org/10.1112/jlms.70425}{doi:10.1112/jlms.70425}.

\bibitem{Lorch1993}
L.~Lorch,
\newblock Some inequalities for the first positive zeros of Bessel
functions,
\newblock \emph{SIAM J. Math. Anal.} \textbf{24} (1993), no.~3, 814--823.
\newblock \href{https://doi.org/10.1137/0524050}{doi:10.1137/0524050}.

\bibitem{DLMF}
NIST Digital Library of Mathematical Functions,
\newblock Chapter 10: Bessel functions, in particular equation (10.9.30),
\S10.21(i), equation (10.21.2), and equation (10.47.3),
\newblock \url{https://dlmf.nist.gov/10}, accessed 16 July 2026.

\end{thebibliography}

\end{document}
